% Cote 119: offprints, reports, courses, CV and letters (typed and printed part): archive cote 119, site page 198 (PDF page 199).
% typescript. Draft compared: none.
% AI-transcribed by Claude Opus 5.5 (claude-opus-5-5, adaptive thinking, effort high), Message Batches API, batch r119_full (msgbatch_01DwZNm5Y7YsKBWyAN4sFs8w), 2026-11-13; not yet checked by eye.
\scanpage{119/198}
\mspage{- 3 -}

(plus généralement, des schémas), qui convienne à l'interprétation
cohomologique des célèbres conjectures de Weil. En effet, la notion tradition-%
nelle d'espace topologique ne suffit pas à traiter le cas des variétés algé-%
briques sur un corps autre que le corps des complexes, la topologie proposée
précédemment par Zariski ne donnant pas lieu à des invariants cohomologiques
``discrets'' raisonnables. A l'heure actuelle, le point de vue des topos, et la
notion de ``localisation'' correspondante, font partie de la pratique quotidienne
du géomètre algébriste, et il commence à se répandre également en
\uline{théorie des catégories} et en \uline{logique mathématique} (avec la démonstration par B.
Lawvere [*] du théorème de Cohen d'indépendance de l'axiome du continu, utili-%
sant une adaptation convenable de la notion de topos). Il n'en est pas encore
de même en topologie et en géométrie différentielle et analytique, malgré cer-%
tains premiers essais dans ce sens (comme la tentative de démonstration par
Sullivan d'une conjecture d'Adams en K-théorie, par réduction à une propriété
de l'opération de Frobenius sur les variétés algébriques en car. $p > 0$).

\uline{Références} : M. Atiyah, F. Hirzebruch, H. Bass, J. Leray, M. Artin, D. Quillen,
M. Karoubi ...

4. \uline{Algèbre} ([15], [16], [18])

Comme l'algèbre homologique, l'algèbre a été pour moi un outil à
développer, et non un but en soi. J'ai parlé au par. 2 de mes contributions à
l'algèbre homologique, et au par. 3 de mes contributions à la K-théorie; celle-%
ci comprend une partie purement algébrique (qui, une fois étendue en une
théorie des $K^{i}$ supérieurs, finira par devenir une partie de l'algèbre homo-%
logique ou homotopique). Ainsi, un certain nombre de mes résultats en géométrie
algébrique se spécialisent en des résultats en algèbre pure, comme la relation
$K(A[t]) \simeq K(A)$ , où $A$ est un anneau. Mises à part ces retombées, on peut
signaler les contributions ci-dessous.

a) \uline{Algèbre catégorique}: En fait, de façon continuelle depuis 1953, je me suis
senti dans l'obligation, au fur et à mesure des besoins, de développer une
panoplie catégorique toujours insuffisante. La plupart des résultats et des
notions ainsi introduites se trouvent développés un peu partout dans [15,16],
notamment dans le premier exposé de SGA 4. Il ne peut être question de passer
en revue ici même sommairement les notions qui sont ainsi entrées dans l'usage
courant. Signalons seulement ici le langage des \uline{univers} (pour éliminer des
difficultés logiques dans la manipulation intensive des catégories), et celui

\margin{198}
